\section{Conclusions}
This research introduces Ecogent, a new multi-agent system designed to work primarily on the user's local computer. The main goal of this system is to solve the major problems of traditional cloud-based AI agents, such as high costs, slow response times, and limited memory. By handling most tasks locally instead of immediately sending them to the cloud, this system proves that complex AI tasks can be completed efficiently and affordably.\\

To make this possible, the system uses a fast local supervisor and a small local AI model to quickly understand what the user wants to do. It also uses a local database (ChromaDB) to build a long-term memory. For example, when the system finishes a web task, it saves the exact steps it took. If the user asks for the same task later, the system simply replays the saved steps for free instead of asking the cloud AI to solve it again. Ultimately, Ecogent provides a smart and sustainable model for future AI agents. It ensures that massive, expensive cloud models are only used for problems that truly require deep reasoning.\\

\section{Future Directions}
While the early version of the Ecogent system proves that a local-first approach works well, there is still important work to be done. A major future goal is to build a smarter routing system that can automatically choose the cheapest and best cloud model for difficult tasks (for example, using a cheaper model like Llama-3 for medium tasks, but saving GPT-4 for the hardest problems).\\

Future work will also add more specialized agents to the system. This includes a Coding Agent with a smart code map for software development, a Data Agent for working with spreadsheets, and an API Agent for connecting to online services. Finally, more research will be done to run large-scale tests comparing Ecogent against traditional cloud agents. These tests will provide exact numbers on how much time, money, and computing power this new system saves in real-world situations.\\

